% Measurement/instrument-calibration methodology short paper -- the
% "instrument-calibration methodology story" identified as a purpose-built
% fit for AIMS/Sci-FM-class venues (see
% ../neurips-ws-2026/VENUE_DECISION.md, "Live/near-live windows" +
% "Per-venue one-liners: Sci-FM/AIMS", and STATE.md's matching PI-scout
% entry: "AIMS/Sci-FM ~ purpose-built for the instrument-methodology
% story"). Dispatched 2026-07-06, drafted 2026-07-08.
%
% SCOPE BOUNDARY (read before editing): this paper reports the MEASUREMENT
% FAILURE MODE and the audit PROCESS that caught it -- a numerical-
% tolerance miscalibration in an eval-time admission gate that produced an
% apparent capacity collapse at d_state=96. It does NOT claim the capacity
% RESULTS themselves (cliff location at d=64, dissolution at d=128,
% coherence-dose-response exoneration) -- those are the companion capacity-
% trilogy paper's own contribution (../workshop-2026/, cited throughout as
% \citet{larson2026capacitycliff}). Numbers describing the d=96/pool-margin
% episode necessarily appear in both papers (both are drawn from the same
% underlying repo record, KEY_ANCHORING_SCALING_DRAFT.md \S15.22-\S15.27,
% and EXPERIMENT_LOG.md's matching entries) -- the CLAIM each paper stakes
% on those numbers differs, per \S1's "Scope and companion paper"
% paragraph below. Never edit workshop-2026/'s files from this paper's
% own drafting session; never edit this paper's files to assert a new
% capacity-boundary result beyond what \S15.22-\S15.27 already record.
%
% VENUE: single-column COLM-2026-style scaffold (COLM's own official
% .sty is not vendored here -- this preamble is deliberately generic
% article-class + geometry so it retargets cleanly to a twocolumn
% NeurIPS-workshop-style scaffold (AIMS/Sci-FM/NeurIPS-ws) by swapping
% the \documentclass line and re-wrapping the abstract, mirroring the
% companion papers' own "reusable scaffold until an official style file
% is confirmed" convention (../workshop-2026/main.tex, ../neurips-ws-2026/
% main.tex).
%
% AUTHOR BLOCK: Sam Larson, affiliation placeholder below -- PENDING PI
% DECISION (see STATE.md "PENDING PI DECISIONS" item 1 and
% ../neurips-ws-2026/VENUE_DECISION.md's venue/author-block items). Do
% NOT fill in \texttt{[affiliation]} without a PI sign-off, matching the
% other two workshop papers' own placeholder convention.
%
% TITLE: working title as given in the task brief, used verbatim below.
%
% PROVENANCE DISCIPLINE: every number in this paper is cited, in a LaTeX
% comment adjacent to its first use, to an exact \S15.NN anchor in
% matrix-thinking/KEY_ANCHORING_SCALING_DRAFT.md and/or the matching
% EXPERIMENT_LOG.md entry, so a verifier can grep and check. No number
% below was invented or estimated; every one was re-read from those two
% files this drafting session.

\documentclass[10pt,letterpaper]{article}

\providecommand{\anon}{0}

\usepackage[margin=1in]{geometry}
\usepackage{times}
\usepackage{microtype}
\usepackage[utf8]{inputenc}
\usepackage[T1]{fontenc}
\usepackage{amsmath,amssymb,amsthm}
\usepackage{graphicx}
\usepackage{booktabs}
\usepackage{array}
\usepackage{xcolor}
\usepackage{xurl}
\usepackage[hidelinks]{hyperref}
\usepackage[numbers,square,sort&compress]{natbib}
\usepackage{caption}
\captionsetup{font=small,labelfont=bf,skip=4pt}

\newcommand{\anonymous}{Anonymous}
\newcommand{\dstate}{d_{\mathrm{state}}}
\newcommand{\niter}{n_{\mathrm{iter}}}

\title{The Cliff That Wasn't: How Numerical Tolerance Miscalibration\\
Manufactured a Capacity Collapse (and What Caught It)}

\if\anon1
  \author{\anonymous{} Authors\\
  \anonymous{} Institution\\
  \texttt{anonymous@example.com}}
\else
  % TODO(PI): affiliation is a placeholder pending sign-off -- do not fill
  % in without approval (see comment block above).
  \author{Sam Larson\\
  \texttt{[affiliation]}\\
  \texttt{email@TBD}}
\fi

\date{}

\begin{document}

\maketitle

\begin{abstract}
A wide-grid scaling sweep on a trained matrix fast-weight memory appeared
to show a capacity collapse: at state dimension $\dstate=96$, 11 of 12 new
load cells ($K \in \{72,78,84,90\}$) failed an eval-time admission gate --
a result that, reported as measured, would have been a headline negative
finding, directly contradicting the same architecture's clean super-linear
capacity growth already measured from $\dstate=64$ to $\dstate=80$. We
report the diagnosis instead. The admission gate's Newton--Schulz
convergence check, tuned against anchor-stabilized \emph{blended} keys,
was miscalibrated for a structurally different, un-stabilized \emph{raw}
key population that one specific evaluation probe pool actually queries. A
purpose-built, four-times independently-audited repro instrument recovered
the true failing tensor and showed every one of 4{,}608 flagged episodes
resolves cleanly under a modestly larger iteration budget ($\niter=20 \to
28$) -- a near-miss, not a structural wall. Recalibrating the admission
\emph{flag} alone (never the underlying capacity measurement) unlocked the
quarantined cells at zero additional GPU spend; the re-fit shows no cliff
through $K/\dstate = 0.9375$. We report four further findings from the
same instrument-audit chain: the recalibrated frontier itself moves with
load (sufficient at $K{=}84$, insufficient at $K{=}90$); the failing
admission leg swaps identity across recalibrations; an archived
exact-ceiling reading did not replicate at a fresh seed; and a diagnosed
confound shows a real, directionally-consistent effect roughly
twelve times its own measured noise floor, correctly withheld from
promotion to a finding because its own diagnostic cell never cleared
admission. We describe the discipline -- pre-registered decision rules
over instrument state, independent adversarial audit rounds with fresh
eyes at each round, negative unit tests exercised (and mutated) to
completion, and verification against raw artifacts rather than summary
logs -- that caught the miscalibration before it became a published
claim, and argue it generalizes to any measurement-heavy empirical ML
claim gated by a numerical admission check.
\end{abstract}

\section{Introduction}
\label{sec:intro}

Empirical claims about learned systems are only as trustworthy as the
instrument that measured them. This is a case study of one instrument
failing quietly, being caught before publication, and being fixed --
reported because the failure mode (a numerical tolerance calibrated
against one data population, silently misapplied to a structurally
different one) and the practice that caught it (pre-registered decision
rules, independent adversarial audit, verification against raw artifacts)
generalize to any measurement-heavy ML claim gated by a numerical
admission check. \textbf{The setup, one sentence:} a trained matrix
fast-weight memory's associative-recall capacity is measured by exact
continuous recovery across a load sweep $K/\dstate$, and a separate,
eval-time admission gate decides which measured cells are trustworthy
enough to enter a capacity fit.

\textbf{Scope and companion paper.} A companion paper
\citep{larson2026capacitycliff} reports this program's capacity
\emph{results}: the cliff located at $\dstate=64$, its dissolution at
$\dstate=128$, and a coherence confound tested and exonerated; its own
Limitations section discloses several of the same raw numbers analyzed
here, as scope-limited caveats on its $\dstate=96$ material. This paper's
claim is different in kind: it is about the measurement \emph{failure
mode itself} -- how a capacity collapse was manufactured by an
uncalibrated instrument, and what process caught it -- not a new
capacity-boundary result. Where a number appears in both papers, this
paper cites it for what the audit trail says about instrument
reliability, not to re-claim it as a located transition.

\section{Setup}
\label{sec:setup}

The testbed is a single-head, single-layer DeltaNet-style fast-weight
memory \citep{schlag2021linear}: a per-entity key-anchoring mechanism
provides a table of $107$ trainable entity keys, and capacity is measured
by held-out compositional recovery at hop depth 4 ($h4$), graded by
\emph{exact continuous} cosine similarity against the true target vector
at threshold $0.9$ -- never argmax over a codebook, the design choice that
keeps a rank or capacity bound from being silently satisfiable by
codebook lookup \citep{nichani2025factual}. $K$ is the number of
independent entity bindings one episode must simultaneously support;
$K/\dstate$ is the capacity ratio the load sweeps below vary.

Not every trained cell's $h4$ is automatically eligible for a capacity
fit: an eval-time \emph{admission gate} must clear four disjoint legs --
a value-salvage tier check, a finite-loss/no-divergence check, a
task-performance floor, and a Newton--Schulz convergence check -- before
a cell enters a fit.\footnote{Exact identifiers, verified against the
live code:
\texttt{value\_salvage\_tier\_pass},
\texttt{finite\_loss\_no\_divergence},
\texttt{task\_performance\_floor\_pass},
\texttt{ns\_converged\_no\_fallback}.}
The leg this paper is about, Newton--Schulz convergence
(\texttt{ns\_converged\_no\_fallback}), gates whether an orthogonalization
procedure -- an ancillary computation \emph{inside the recovery-probe
scoring path}, not part of the trained model itself -- converges to a
residual under $\mathrm{resid\_tol}=0.01$ within a fixed iteration budget
$\niter$ (production setting $\niter=20$).
% [PROVENANCE] Four-leg admission definition, resid_tol=0.01, production
% niter=20: KEY_ANCHORING_SCALING_DRAFT.md \S15.23 ("cross-checked
% directly against compute_geo3_admission" table + "resid_tol=0.01,
% exactly key_anchoring.GATE2_RESID_TOL").

\section{The Phenomenon: An Apparent Capacity Collapse at $\dstate=96$}
\label{sec:phenomenon}

A wide-grid load sweep at $\dstate=96$ ran 12 new cells
($K\in\{72,78,84,90\}$, 3 seeds each) alongside 3 reused $K{=}69$ cells.
% [PROVENANCE] Wave design/scope: KEY_ANCHORING_SCALING_DRAFT.md \S15.22
% ("12 new d=96 cells... reusing the 3 already-archived K=69/d=96 cells").
Training was clean at every single cell -- $h1$ (the in-distribution
training-health guard) read exactly $1.0$ at all 19, and
$n\_geo3\_fallback\_train\_steps=0$ everywhere -- but of the 12 new
cells, only \textbf{1 was admissible}: $K{=}78$, $K{=}84$, and $K{=}90$
each failed admission on all 3 of their own seeds (0/3 each), an overall
new-cell admissibility rate of \textbf{5/16 (31.25\%)}.
% [PROVENANCE] Per-cell admissibility table, h1=1.0 guard, training-health
% clean: KEY_ANCHORING_SCALING_DRAFT.md \S15.22 per-cell table + "h1 guard
% ... reads exactly 1.0 at all 19 cells" + "Admissibility, new cells only:
% 5/16 (31.25%)".
Reported as measured, this is the shape of a genuine capacity cliff: a
sharp admissibility collapse rising from $0/30$ at $K/\dstate \le 0.71875$
to $3/3$ inadmissible at every $K/\dstate \ge 0.8125$ -- directly
contradicting the clean, non-degenerate super-linear growth already
measured between $\dstate=64$ and $\dstate=80$
($x_0 = 0.5455 \to 0.6779$, \citealp{larson2026capacitycliff}): a
newsworthy, and as it turned out wrong, negative result.
% [PROVENANCE] Admissibility-collapse rate by K/d:
% KEY_ANCHORING_SCALING_DRAFT.md \S15.22 addendum ("intensifies sharply
% above K/d~=0.75: 0/30 at ... K/d<=0.71875, 1/3 at K/d=0.75 ... 3/3 at
% each of K/d=0.8125/0.875/0.9375").
% [PROVENANCE] x0=0.5455 (d=64) is the companion paper's own headline
% result (cited to it, not re-derived here); x0=0.6779 (d=80) is this
% program's own escalated (n=5 seed) re-fit, established within the cited
% range: KEY_ANCHORING_SCALING_DRAFT.md \S15.22 ("sigmoid fit: x0=0.6779
% w=0.0479 L=0.9994 rss=0.00024").

The \emph{first} mechanistic hypothesis was itself wrong, and catching
that was the episode's first lesson. The initial writeup attributed the
failure to the trained anchor table's own Newton--Schulz convergence at
high $K/\dstate$. An independent follow-up diagnostic tested that claim
directly on six pulled checkpoints (four failing, two passing controls)
and refuted it: the anchor table converges cleanly at $\niter=20$ with
residuals \textbf{$\sim$7{,}000--8{,}000$\times$ below} the $0.01$
tolerance, indistinguishable between failing and passing cells. The real
failure traced to a single evaluation probe pool
(\texttt{C17\_heldout\_entities}) that is, by construction, architecturally
anchor-\emph{bypassed} -- the anchor table's own conditioning cannot input
a computation it never participates in. The original mechanistic claim
was formally retracted before any fix was attempted.
% [PROVENANCE] Anchor-table residual margin, C17-exclusive/anchor-bypass
% finding, retraction: KEY_ANCHORING_SCALING_DRAFT.md \S15.23 VERDICT
% MISDIAGNOSED-ARTIFACT ("every measured residual sits ~7,000x-8,000x
% below the admission threshold" + "the failure is 100%
% C17_heldout_entities-exclusive... \S15.22's own mechanistic-root-cause
% claim is retracted").

\section{The Catch: Instrument Saturation, Not Model Failure}
\label{sec:catch}

With the wrong object identified and ruled out, a purpose-built repro
instrument was designed to recover the \emph{true} failing tensor: C17's
raw, pre-Newton--Schulz post-conv keys for held-out entities, which the
original wave's checkpoint writer never saved. The design went through
\textbf{four independent adversarial attack rounds} before any GPU cell
launched, each catching real defects, not merely stylistic ones.
% [PROVENANCE] Repro instrument design, Rev 0 -> Rev 4 attack-round
% history: KEY_ANCHORING_SCALING_DRAFT.md \S15.24 (status notes at each
% revision).
Round 1's FATAL fix closed a gap where an empty telemetry dump (a
plausible outcome, given this hardware's own measured fixed-seed
run-to-run nondeterminism) would let a universally quantified check
vacuously pass and silently emit an unearned \textsc{tolerance-
miscalibration} verdict -- in the design's own words, ``an unearned
`unlock the 11 cells' paper claim.''
% [PROVENANCE] Round-1 FATAL-1 (zero/low-event guard):
% KEY_ANCHORING_SCALING_DRAFT.md \S15.24 fix-map table, FATAL-1 row
% (quoted verbatim in-text).
Round 4's MAJOR fix closed a ``two-seeds trap'': the offline analysis
reconstructs the training/held-out entity partition from a
\emph{hardcoded} seed decoupled from the cell's own launch seed --
naively threading the launch seed in would silently reconstruct the
wrong partition and produce a confidently wrong verdict on healthy code.
% [PROVENANCE] Round-4 MAJOR-1 ("two seeds trap"):
% KEY_ANCHORING_SCALING_DRAFT.md \S15.24 Rev 4 status note.

The instrument's own decision rule was pre-registered as a strict
precedence order over telemetry state, walked mechanically at harvest
time: a pool-membership precheck (any single violation is dispositive,
structural, no statistics required) $\to$ a reproduction-count floor
$\to$ live/offline Newton--Schulz agreement $\to$ an $\niter$ sweep. Every
step ran, in order, on the real cell ($K{=}84$, seed$=1940$): the
reconstruction gate passed ($107/107$ set-equal); \textbf{zero}
pool-membership violations across all 36 dumped events (ruling out an
instrument bug); the reproduction floor cleared at $12\times$ its own
minimum; \textbf{zero} live/offline disagreements across all
\textbf{4{,}608} examined episodes (ruling out an instrument bug a
second, independent way); and the $\niter$ sweep resolved
\textbf{every single one} of those 4{,}608 episodes by $\niter\le28$
($4{,}313$ at $\niter=24$, the remaining $295$ at $\niter=28$, zero left
unresolved even at $\niter=32$ or $40$) $\to$ \textsc{tolerance-
miscalibration}.
% [PROVENANCE] Full decision-rule walk, per-step results, 4,608-episode
% n_iter-sweep tally: KEY_ANCHORING_SCALING_DRAFT.md \S15.25.1 (also
% independently re-tallied in EXPERIMENT_LOG.md's matching 2026-07-08
% C17 repro run+harvest entry).

The magnitude makes the "near-miss, not a wall" reading concrete: at the
final checkpoint, the worst-case flagged episode's residuals ranged
\textbf{0.696--1.371} against the $0.01$ tolerance -- 70$\times$--
137$\times$ over threshold -- and this near-miss population
\emph{worsened monotonically} over training (13--27 of 128 anomalous rows
at step 16{,}000, all 128/128 by step 20{,}000). Yet it still cleared
cleanly by $\niter=28$. A structural wall would be expected to keep
worsening past $\niter=40$; a miscalibrated tolerance resolves at a
fixed, modest iteration bump, and did.
% [PROVENANCE] Event-0 residual range, monotonic worsening over training:
% KEY_ANCHORING_SCALING_DRAFT.md \S15.25.1 ("event-0 range 0.696-1.371"
% + "checkpoint 16000's own 12 events show only 13-27/128 anomalous rows
% ... checkpoints 18000 and 20000 show ALL 128/128").

Because the underlying capacity measurement was never in question --
training converged cleanly at all 11 quarantined cells regardless of
admission status -- only the admission \emph{flag}, never $h4$, was
recalibrated, and only for the 11 cells independently verified to share
the identical failure signature (one pre-existing, out-of-scope anomaly
at $K{=}69$ was deliberately \emph{not} flipped, a disclosed scope
decision, not an oversight). This unlocked the full grid at
\textbf{zero additional GPU spend}. The re-fit found no cliff anywhere
through $K/\dstate=0.9375$: per-$K$ mean $h4$ reads
$0.9592 (K{=}69), 0.9216 (K{=}72), 0.9326 (K{=}78), 0.9581 (K{=}84),
1.0000 (K{=}90)$ -- non-monotonic, but nowhere collapsing.
% [PROVENANCE] Unlock scope (11 of 12 cells, K=69 excluded), zero-GPU
% mechanism, re-fit per-K means: KEY_ANCHORING_SCALING_DRAFT.md \S15.25.4
% ("THE UNLOCK") and \S15.25.5 (per-K table).

\section{General Findings: Admission Checks Are Instrument-Relative}
\label{sec:findings}

The follow-up diagnostic chain (a zero-GPU power check, then a two-cell
control diagnostic at the recalibrated $\niter=28$) surfaced five findings
that generalize beyond this one instrument.

\textbf{(a) The $\niter$-sufficiency frontier moves with load.}
$\niter=28$ resolved every flagged episode at $K{=}84$ (\S\ref{sec:catch}).
A follow-up control cell at $K{=}90$, run under the \emph{same}
recalibrated $\niter=28$, stayed fallback-inadmissible: a frontier
measured at one $K/\dstate$ does not transfer to a higher one without
re-verification.
% [PROVENANCE] K=90 fallback-inadmissible at n_iter=28:
% KEY_ANCHORING_SCALING_DRAFT.md \S15.27 escalated finding (2).

\textbf{(b) Admission legs swap failure modes across recalibrations.}
At $\niter=28$, $K{=}84$ fails a \emph{different} leg
(\texttt{value\_salvage\_tier\_pass}, $0.09307 < 0.10$, training
convergence now fully clean) than the original $\niter=20$ population,
which failed convergence while value-salvage \emph{passed}
($0.102$--$0.126$). Fixing one failure mode can silently expose the next.
% [PROVENANCE] Leg-swap at K=84: KEY_ANCHORING_SCALING_DRAFT.md \S15.27
% escalated finding (1).

\textbf{(c) An archived exact ceiling did not replicate.} $K{=}90$'s
original 3 seeds all read exact $h4=1.0000$; a fresh seed at the
recalibrated $\niter=28$ read $0.9725$ -- ``exact'' described the specific
seeds measured, not a property of the load.
% [PROVENANCE] K=90 fresh h4=0.972526: KEY_ANCHORING_SCALING_DRAFT.md
% \S15.27 descriptive finding (3) / EXPERIMENT_LOG.md matching entry.

\textbf{(d) A diagnosed confound is real in direction but was correctly
withheld.} Restricting $K{=}84$'s eval-time entity-pool margin to match
$K{=}90$'s own overlap fraction raised $h4$ by $+0.0330$ -- roughly
\textbf{12$\times$} the measured eval-sampling noise floor ($0.0028$, the
larger of two independent noise draws). The mechanism is directionally
real, but because \emph{both} of this diagnostic's own cells stayed
inadmissible even at $\niter=28$, the pre-registered routing correctly
refused to promote it to a finding -- the discipline that caught the
original miscalibration also declined to over-claim its own follow-up.
% [PROVENANCE] Pool-restriction shift +0.0330 vs noise floor 0.0028
% (~12x), DEGENERATE_CELL routing: KEY_ANCHORING_SCALING_DRAFT.md \S15.27
% ("K=84's pool-restriction shift = +0.0330 vs a measured noise floor of
% 0.0028 ... ~12x noise") + EXPERIMENT_LOG.md matching harvest entry
% (0.033040 / 0.002767).

\textbf{(e) Structural checks need exact thresholds, never tolerance
slack borrowed from a numerical context.} This instrument's own attack
history supplies the counter-example: an early draft of the
pool-membership precheck reused a $\ge2$-event recurrence floor built for
a genuinely noisy \emph{numerical} disagreement check, applied unexamined
to a \emph{structural} fact -- pool membership has no floating-point
component; an id either is or is not in a disjoint set. Under the reused
floor, a single real violation in a small event set would be silently
excluded, and the verdict would continue, confidently and wrongly, toward
\textsc{tolerance-miscalibration}. The fix split the floor: a single
structural violation is dispositive on its own; the recurrence floor is
retained only for the numerical check it was designed for.
% [PROVENANCE] Structural-vs-numerical floor split (the exact-threshold
% lesson, this instrument's own instance of it):
% KEY_ANCHORING_SCALING_DRAFT.md \S15.24 fix-map, FATAL-1 row (quoted
% and paraphrased in-text). Generalizes the same rule already recorded
% project-wide in this repo's CLAUDE.md hard-rules list.

\begin{table}[t]
\centering
\footnotesize
\renewcommand{\arraystretch}{0.95}
\begin{tabular}{@{}p{0.6cm}p{3.3cm}p{3.6cm}@{}}
\toprule
\S & Wave & Verdict \\
\midrule
15.22 & $\dstate{=}96$ wide-grid load sweep & AMBIGUOUS (data-quality collapse) \\
15.23 & anchor-table NS diagnostic & MISDIAGNOSED-ARTIFACT \\
15.24 & C17 repro instrument, design & DESIGN-CLEARED (4 audit rounds) \\
15.25 & C17 repro instrument, harvest & TOLERANCE-MISCALIBRATION $\to$ unlock $\to$ re-fit AMBIGUOUS (non-monotonic, no cliff) \\
15.26 & scatter-resolution power check & finding registered (zero GPU); 10-cell grid killed \\
15.27 & K90 pool-margin control & DEGENERATE\_CELL (3 escalated findings, PI-gated) \\
\bottomrule
\end{tabular}
\caption{The full diagnostic chain and its verdicts, each anchored to its
own \S in \texttt{KEY\_ANCHORING\_SCALING\_DRAFT.md}. No step skipped or
silently absorbed a prior step's finding.}
\label{tab:chain}
\end{table}

\section{The Methodology That Caught It}
\label{sec:methodology}

Table~\ref{tab:chain} names six waves; four independent adversarial audit
rounds landed on the repro instrument's own design alone
(\S\ref{sec:catch}), and a fifth wave's design took two further rounds
plus a dedicated verification pass before any GPU cell launched.
% [PROVENANCE] Pool-margin wave's own 2-round-plus-verify history:
% KEY_ANCHORING_SCALING_DRAFT.md \S15.26.9/\S15.26.10 (Rev 1/Rev 2
% fix-maps) + the "\S15.26 ROUND-3 VERIFY PASS" block (commit c66b3f6,
% verdict DESIGN-CLEARED-FOR-BUILD).
That later rounds keep finding real defects -- round 4's two-seeds trap
(\S\ref{sec:catch}) is not a first-round-only phenomenon -- is itself the
argument against stopping at one audit pass.

This one case study sits inside a well-documented broader pattern:
reported ML results are known to be sensitive to analytic and
implementation choices invisible in the final number
\citep{silberzahn2018many, damour2020underspecification}, that empirical
rigor has not obviously kept pace with the field's own pace of progress
\citep{sculley2018winners, lipton2018troubling}, and that unaccounted
sources of variance routinely change which of two systems looks better
\citep{henderson2018deep, bouthillier2021accounting}; reporting-standard
proposals such as REFORMS \citep{kapoor2023reforms} formalize much of the
raw-artifact verification and pre-registration this instrument-audit
chain practiced informally, ahead of any such checklist.

Four practices recur across every wave in Table~\ref{tab:chain} and are
offered here as transferable, not architecture-specific:

\begin{itemize}
  \item \textbf{Pre-registration of outcome semantics.} The repro
    instrument's six-step decision rule mapped every possible telemetry
    state to exactly one named verdict, fixed \emph{before} the cell that
    would decide this paper's own headline claim ran.
  \item \textbf{Negative tests exercised, and mutated, to completion.} A
    synthetic empty-sink fixture had to emit \textsc{no-repro}, never a
    verdict; a synthetic single-violation pool-mismatch fixture had to
    emit \textsc{instrument-bug} unconditionally; the pool-margin wave's
    own safety guards were \emph{mutation-tested} -- reverting a guard had
    to make its own regression suite fail, not merely pass when correct.
% [PROVENANCE] Negative-test fixtures (E1/E4 pool-mismatch cases, empty-
% sink NO-REPRO): KEY_ANCHORING_SCALING_DRAFT.md \S15.24 fix-map
% (FATAL-1/MAJOR-3 rows, Stage -1 items). Mutation-testing of pool-margin
% guards: \S15.26.10 "BUILD AUDIT" block ("verified by re-execution AND
% by mutation-testing the new guards themselves").
  \item \textbf{Spend compute only where it changes the answer.} A
    registered 10-cell, $4.27$--$8.54$ GPU-h seed-escalation sweep was
    killed and replaced by a $360{,}000$-trial, zero-GPU bootstrap power
    check (8 seeds, an independent reimplementation, a positive control
    proving the check has teeth) that answered the same question at CPU
    cost.
% [PROVENANCE] Killed 10-cell grid, 360,000-trial power check, positive
% control: KEY_ANCHORING_SCALING_DRAFT.md \S15.26.1.
  \item \textbf{Verify against raw artifacts, not summary logs.} Every
    headline number here -- the admissibility table, the 4{,}608-episode
    $\niter$-sweep tally, the pool-margin shift and noise floor -- was
    independently re-derived, at the harvest that reported it, from raw
    JSON/\texttt{.pt} dumps or a fresh read-only box pull, not re-read
    from an earlier commit message.
% [PROVENANCE] "Independently re-verified... not merely re-asserted from
% the commit message" discipline, applied at every harvest in Table 1:
% KEY_ANCHORING_SCALING_DRAFT.md \S15.22, \S15.25.1, \S15.27 (each
% harvest's own opening verification paragraph).
\end{itemize}

The economics are worth stating plainly: the repro instrument that
identified \textsc{tolerance-miscalibration} cost \textbf{0.82 GPU-h};
the follow-up pool-margin diagnostic cost \textbf{1.09 GPU-h} -- under
\textbf{2 GPU-h combined} to interrogate an instrument whose uncorrected
reading would have discredited \textbf{6.33 GPU-h} of clean training
compute already spent, and shipped as a false 11-of-12-cell capacity
collapse. Auditing the instrument was cheap relative to the compute it
was protecting.
% [PROVENANCE] Repro-instrument realized cost 0.82 GPU-h:
% KEY_ANCHORING_SCALING_DRAFT.md \S15.25.3. Pool-margin diagnostic
% realized cost 1.0933 GPU-h: \S15.27 / EXPERIMENT_LOG.md matching entry.
% Original wide-grid wave's own realized cost 6.3331 GPU-h: \S15.22
% "Realized GPU-h vs. ceiling" table.

\section{Limitations}
\label{sec:limitations}

\textbf{Single architecture family.} All findings come from one
DeltaNet-style fast-weight testbed and one specific admission gate
(Newton--Schulz residual tolerance inside an eval-time recovery probe).
No claim is made that this exact failure mode, or this exact fix,
recurs in other admission gates, architectures, or codebases -- only that
this one, checked, was miscalibrated and is now fixed.

\textbf{The ceiling fine structure remains genuinely open.} Finding (d)'s
pool-margin mechanism is real in direction but the diagnostic built to
confirm or refute it landed \textsc{degenerate\_cell}: neither of its own
two cells cleared admission even at $\niter=28$, so its own adjudication
triggers never fired. This is registered for a follow-up design round and
is explicitly \textbf{PI-gated}, not run as of this writing.
% [PROVENANCE] DEGENERATE_CELL routing, PI-gated §15.28-class follow-up:
% KEY_ANCHORING_SCALING_DRAFT.md \S15.27 ("Next step: ... registered for
% the PI check-in and a \S15.28-class design round; NOT self-launched")
% + STATE.md "PENDING PI DECISIONS" item 3.

\textbf{One instrument, one recalibration, no counterfactual.} This paper
documents one class of numerical miscalibration (Newton--Schulz residual
tolerance, calibrated against a blended-key population, silently
misapplied to a raw-key one) and one recalibration of it -- not a general
certificate that every admission gate in this codebase, or any other, is
now trustworthy. The methodology narrative is reconstructed from this
project's own audit trail, not a controlled comparison against a matched
project that skipped the audit steps; we cannot report a rate of
miscalibration a lighter process would have missed, only that this
process caught the one it did.

\textbf{Companion-paper boundary, restated.} The capacity numbers this
paper cites for provenance ($x_0$ progression, the $\dstate=96$ re-fit,
the pool-margin diagnostic's own descriptive readings) are also reported,
as scope-limited caveats, in \citet{larson2026capacitycliff}. That paper's
claim is about the located transition and its dissolution with scale;
this paper's claim is about the measurement failure mode and the process
that caught it before it became one.

\bibliographystyle{plainnat}
\bibliography{refs}

\end{document}
